\session{5}{10/17}{2限}{2}{業務プロセスとAIによる自動化}

\goals{
  \item 業務プロセスを可視化する意味と方法（As-Is／To-Be）を理解する。
  \item RPA、生成AI、AIエージェントの違いと、自動化の進展を説明できる。
  \item 何を自動化し、何を人が担うかを判断する基準を持つ。
}

\guidesec{解説}

\topic{業務プロセスの見える化}
業務プロセスとは、ある目的（受注を処理する、新入社員を受け入れる）を達成するための一連の作業の流れです。
多くの職場では、プロセスは担当者の頭の中にあり、文書化されていません。自動化や改善の前に、
まず現状の流れ（As-Is）を可視化することが出発点です。これを体系的に行う経営手法を BPM（ビジネスプロセス管理）と呼び、
可視化→分析→改善（To-Be の設計）→実行→監視の循環で進めます。

可視化には、作業の順序を矢印でつなぐフローチャートや、部門ごとのレーン（泳路）で担当を分けて描くスイムレーン図が使われます。
可視化の際に確認する点は次の通りです。
\begin{itemize}
  \item 各作業の入力（何を受け取るか）と出力（何を次に渡すか）。
  \item 作業の担当者と、判断が必要な分岐点（承認、例外処理）。
  \item 作業ごとの所要時間、件数、待ち時間。待ち時間が全体の時間の大半を占めることが多い。
  \item 手作業の転記や二重入力など、情報システムの「つなぎ目」の無駄。
\end{itemize}
可視化せずに自動化すると、無駄な手順をそのまま高速化するだけになります。

\topic{RPA から生成AI、AIエージェントへ}
業務の自動化は、対象とする作業の性質によって段階的に広がってきました。
\par\noindent\begin{gtabh}{colspec={Q[l,wd=24mm] X[l] X[l] X[l]}}
  技術 & 得意な作業 & 例 & 限界 \\
  RPA & 手順が固定された定型的なPC操作の反復 & 受注メールの内容を基幹システムに転記、複数システムからの集計 & 画面や様式が変わると止まる。判断はできない。 \\
  生成AI & 非定型の文章の作成・要約・分類、判断の支援 & 問い合わせへの回答案の作成、議事録の要約、契約書の要点抽出 & 正確さの保証がない。人の確認が前提。 \\
  AIエージェント & 目標を与えると計画を立て、ツールを使って複数の手順を自律的に実行 & 「来週の出張を手配して」で検索・予約・報告まで行う & 誤った行動が自動で実行される危険。権限と監視の設計が必要。 \\
\end{gtabh}
これらは置き換えの関係ではなく、組み合わせて使われます。たとえば、生成AIが問い合わせメールを分類して回答案を作り、
RPA がその結果を顧客管理システムに登録する、という構成です。AIエージェントはこの組み合わせをAI自身が行う段階と言えます。

\topic{何を自動化し、何を人が担うか}
自動化の対象を選ぶ基準を2つの軸で整理します。
\begin{itemize}
  \item \textbf{定型性}　手順とルールが明確で、例外が少ないか。
  \item \textbf{量}　件数が多く、繰り返されるか。
\end{itemize}
定型で量が多い業務は、自動化の効果が最も大きく、最初に取り組む対象です。定型だが量が少ない業務は、
自動化の費用に見合わないことがあります。非定型で量が多い業務（問い合わせ対応など）は、生成AIによって自動化の候補に入ってきましたが、
人の確認を組み込む設計が必要です。非定型で量が少ない業務（経営判断、例外対応）は人が担います。

さらに、次の点も考慮します。
\begin{itemize}
  \item \textbf{誤りの影響}　誤った処理が顧客や法令に関わる場合は、自動化しても人の確認を残す。
  \item \textbf{人の役割の変化}　自動化後に人が担う作業（例外処理、品質の確認、顧客との関係）を設計する。
  \item \textbf{プロセスの改善が先}　自動化の前に、不要な手順や承認をなくす。
\end{itemize}

\topic{キーワード}
\par\noindent\begin{keywords}{}
  BPM & 業務プロセスを可視化・分析・改善・実行・監視する循環で管理する経営手法。 \\
  As-Is／To-Be & 現状のプロセスと、改善後のあるべきプロセス。 \\
  RPA & 定型的なPC操作をソフトウェアロボットで自動化する技術。 \\
  AIエージェント & 目標に向けて自律的に計画し、ツールを使って複数の手順を実行するAI。 \\
  定型性と量 & 自動化の対象を選ぶ2つの軸。定型で量が多い業務から取り組む。 \\
  Human-in-the-loop & 自動化された処理に人の確認・判断を組み込む設計（第11回で詳述）。 \\
\end{keywords}

\guidesec{演習課題}

\exercise{1}{AI演習：日常の業務を手順に分解し、AIに任せられる段階を分析する}
\instr{自分が経験した業務（アルバイトの開店準備、サークルのイベント運営、レポートの作成など）を1つ選び、次のプロンプトでAIと一緒に手順に分解してください。分解した各手順について、「RPA向き」「生成AI向き」「AIエージェント向き」「人が担う」のいずれかに分類し、理由を書いてください。}
\begin{promptbox}[業務の分解]
私は「（業務の内容）」を担当しています。この業務を、開始から完了までの手順に分解してください。各手順について、入力（受け取るもの）、出力（次に渡すもの）、判断が必要かどうか、繰り返しの回数を表にしてください。分からない点は私に質問してください。
\end{promptbox}

\exercise{2}{自動化の優先順位}
\instr{次の5つの業務を「定型性」と「量」の2軸で位置づけ、自動化の優先順位をつけてください。順位の理由と、自動化後も人が担うべき部分を書きます。}
\begin{itemize}
  \item (A) 毎日300件届く注文メールの内容を受注システムに入力する
  \item (B) 月に1回、各部門の経費の集計表を作る
  \item (C) 1日50件の顧客からの問い合わせに回答する
  \item (D) 年に数件の大口顧客との契約条件を交渉する
  \item (E) 毎週、SNS に投稿する販促文を作る
\end{itemize}

\exercise{3}{AIエージェントに任せる前に決めること}
\instr{「取引先からの請求書をメールで受け取り、内容を確認して会計システムに登録し、支払いを予約する」という業務をAIエージェントに任せる場合、事前に決めておくべきこと（権限、確認の方法、例外の扱いなど）を5つ挙げてください。}

\guidesec{模範解答}

\begin{answer}{演習1}
例：サークルのイベント（新歓）運営
\par\noindent\begin{gtabh}{colspec={Q[l,wd=30mm] X[l] Q[l,wd=24mm] X[l]}}
  手順 & 入力→出力 & 分類 & 理由 \\
  日程と会場の候補出し & メンバーの空き状況→候補3案 & 生成AI向き & 条件から案を複数出す作業。最終決定は人。 \\
  会場の予約 & 候補→予約確定 & AIエージェント向き & Web で空きを調べて申請する一連の手順。ただし支払いは人が確認。 \\
  告知文の作成 & 日程・内容→SNS 投稿文 & 生成AI向き & 下書きを出させ、人が修正する。 \\
  参加申込の集計 & フォームの回答→名簿 & RPA向き & 定型の転記と集計。 \\
  当日の受付と対応 & 来場者→名簿の確認 & 人が担う & 例外（飛び入り参加、トラブル）が多く、関係づくりの場。 \\
  振り返りの議事録 & 録音→要約 & 生成AI向き & 要約は得意。決定事項は人が確認。 \\
\end{gtabh}
AIと対話して気づく点：自分では1つの作業だと思っていた「告知」が、文案作成・画像作成・投稿・返信対応に分かれ、
向き不向きが異なること。「人が担う」部分は、関係づくりと例外対応に集中する。
\end{answer}

\begin{answer}{演習2}
\par\noindent\begin{gtabh}{colspec={Q[c,wd=10mm] Q[l,wd=14mm] Q[l,wd=14mm] Q[c,wd=10mm] X[l]}}
  業務 & 定型性 & 量 & 順位 & 理由と人が担う部分 \\
  A & 高い & 多い & 1 & RPA（または生成AIで内容抽出＋RPAで入力）の典型。様式が崩れたメールは人が処理。 \\
  C & 中 & 多い & 2 & 生成AIで回答案を作り、定型の問い合わせは自動返信、それ以外は人が確認して送る。 \\
  E & 中 & 中 & 3 & 生成AIで下書き。ブランドの表現や炎上の判断は人。 \\
  B & 高い & 少ない & 4 & 定型だが月1回。表計算の自動化で十分で、RPA の投資は見合わないことが多い。 \\
  D & 低い & 少ない & 5 & 交渉は人が担う。AIは過去の条件の整理や想定問答の準備に使う。 \\
\end{gtabh}
\end{answer}

\begin{answer}{演習3}
\begin{enumerate}[label=\arabic*.,leftmargin=2em]
  \item 権限の範囲：エージェントが登録・予約できる金額の上限（例：10万円まで）。それ以上は人の承認を必須にする。
  \item 確認の方法：登録前に、請求書の内容と登録内容の対照表を担当者に提示し、承認を得てから実行する。
  \item 例外の扱い：発注記録と一致しない請求書、初めての取引先、様式が読み取れない請求書は人に回す。
  \item 記録：エージェントが行ったすべての操作と判断の理由をログに残し、監査できるようにする。
  \item 停止と責任：誤った支払いが起きたときの取り消し手順と、最終責任者（経理責任者）を決めておく。
\end{enumerate}
自律性が高いほど、「何をしてよいか」より「何をしてはいけないか」と「誰が止めるか」の設計が重要になります。
\end{answer}

\guidesec{出席確認クイズの解説}
講義サイトの出席確認ページ（第5回）の5問です。
% 出席確認クイズ 第05回　業務プロセスとAIによる自動化
%   attendance/attendance05.html から生成（解説は本ガイド用に加筆）。

\quizq{1}{RPAの説明として適切なものはどれでしょうか?}%
  {ネットワーク機器の一種}%
  {工場で働く物理的な産業用ロボット}%
  {\correct{定型的なPC操作をソフトウェアロボットで自動化するしくみ}}%
  {会計基準の名称}%
  {正解は (c)。RPA（Robotic Process Automation）は、パソコン上の定型作業（転記、集計、入力）をソフトウェアロボットで自動化する技術です。(b) の物理的な産業用ロボットとは別物です。ルールが明確で手順が変わらない作業に向きます。}

\quizq{2}{RPAと生成AIの違いとして適切なものはどれでしょうか?}%
  {RPAは自然な文章を生成できる}%
  {\correct{RPAは決まった手順の反復に強く、生成AIは非定型の文章生成や判断支援に向く}}%
  {生成AIは定型処理しかできない}%
  {両者はまったく同じもの}%
  {正解は (b)。RPA はルール化された定型作業の反復に強く、生成AIは文章作成・要約・判断支援など非定型の作業に向いています。両者は競合ではなく、RPA が生成AIの出力を基幹システムに入力するといった組み合わせで使われます。(a)(c)(d) は特徴を取り違えています。}

\quizq{3}{「AIエージェント」の説明として適切なものはどれでしょうか?}%
  {人間の営業担当者のこと}%
  {一つの質問に一度だけ答えるしくみ}%
  {AIの利用契約書のこと}%
  {\correct{目標に向けて自律的に計画し、ツールを使って複数の手順を実行するAI}}%
  {正解は (d)。AIエージェントは、目標を与えられると自ら計画を立て、検索やシステム操作などのツールを使って複数の手順を実行し、結果を確認しながら進めるAIです。(b) はチャットの一問一答で、エージェントではありません。自律性が高い分、第13回で扱うガバナンスがより重要になります。}

\quizq{4}{BPM(ビジネスプロセス管理)で最初に行うこととして適切なものはどれでしょうか?}%
  {担当者を減らす}%
  {すべての業務を即座に自動化する}%
  {\correct{業務の流れを可視化する}}%
  {新しいシステムを購入する}%
  {正解は (c)。BPM（ビジネスプロセス管理）は、現状（As-Is）の可視化から始めて、分析・改善（To-Be）・実行・監視を繰り返します。可視化をせずに自動化すると、無駄な手順を高速化するだけになります。(a)(b)(d) は手段を目的と取り違えた例です。}

\quizq{5}{自動化の対象を選ぶ際の考え方として適切なものはどれでしょうか?}%
  {すべての業務を無条件に自動化する}%
  {\correct{判断が曖昧で例外の多い業務は人が担い、定型で量の多い業務から自動化する}}%
  {例外の多い業務から優先して自動化する}%
  {自動化は一切行わない}%
  {正解は (b)。自動化の効果は「定型性 × 量」で決まります。定型で件数の多い業務から自動化し、判断が曖昧で例外の多い業務は人が担うのが基本です。(c) は優先順位が逆で、(a)(d) は極端です。}

